\documentclass[11pt]{article}
\usepackage[margin=1in]{geometry}
\usepackage[T1]{fontenc}
\usepackage{lmodern}
\usepackage{amsmath,amssymb,amsthm}
\usepackage{booktabs}
\usepackage{array}
\usepackage{listings}
\usepackage{xcolor}
\usepackage[hidelinks]{hyperref}

\newtheorem{theorem}{Theorem}
\newtheorem{lemma}[theorem]{Lemma}
\newtheorem{proposition}[theorem]{Proposition}
\newtheorem{corollary}[theorem]{Corollary}
\theoremstyle{definition}
\newtheorem{definition}[theorem]{Definition}
\theoremstyle{remark}
\newtheorem{remark}[theorem]{Remark}

% head marker: state s reading symbol x
\newcommand{\hd}[2]{[\mathrm{#1}\,#2]}
\newcommand{\Z}{0^{\infty}}
\newcommand{\coq}[1]{\texttt{#1}}
\newcommand{\W}{\mathcal{W}}

\lstdefinelanguage{Coq}{
  morekeywords={Theorem,Lemma,Proof,Qed,Definition,Fixpoint,Inductive,match,with,end,exists,forall,fun,let,in,Some,None},
  sensitive=true, morecomment=[s]{(*}{*)}}
\lstset{language=Coq,basicstyle=\small\ttfamily,keywordstyle=\bfseries,
  columns=fullflexible,frame=single,framerule=0.3pt,xleftmargin=1em,xrightmargin=1em}

\title{A Proved Transcript Grammar for the BB(5) Champion}
\author{Vishnu Vaidyanathan\\
\small \texttt{vishnu@edgi.in}\\
\small ORCID 0009-0001-6244-5262\\
\small \url{https://github.com/VishnuVaidyanathan/BB5-decimal-macrogrammar}}
\date{November 2026}

\begin{document}
\maketitle

\begin{abstract}
The Marxen--Buntrock machine is the 5-state Busy Beaver winner: it halts after
$47{,}176{,}870$ steps. Its high-level behaviour has been known since Buro (1990) and
Michel (1993), who showed that it follows a Collatz-like map. We study the same machine
through its \emph{transcript}, the sequence of state--symbol pairs it visits, written
as decimal digits. We give a closed-form grammar for the entire transcript: it is
\texttt{0246} followed by 15 layers, each built from a repeated word $(159)^M$, an
opening gate (\texttt{02} or \texttt{1}), a ladder of rungs
$4\,7^{n+2}\,6\,0\,3^{n+3}\,2$, a final rung $4\,7^{n+2}\,6\,0\,3^{n+4}\,2$ and the
closing word \texttt{46}; the last layer ends in \texttt{158}. We prove in Coq, with no
axioms, that the transcript of the machine from the blank tape is exactly this
grammar, and that our layer map is the Buro--Michel map. The step count and the number
of ones on the final tape follow from the proved configurations by short exact
calculations.
\end{abstract}

% ------------------------------------------------------------------
\section{Introduction}

The Busy Beaver value $S(5)=47{,}176{,}870$ was determined in 2024 by the bbchallenge
collaboration, with a proof checked in Coq~\cite{bbchallenge,coqbb5}. The value is set
by the machine found by Marxen and Buntrock in 1989~\cite{marxen}:
\[
\texttt{1RB1LC\_1RC1RB\_1RD0LE\_1LA1LD\_1RZ0LA}.
\]
Its long run is well understood at a high level. Buro~\cite{buro} and, independently,
Michel~\cite{michel93} showed that it follows a Collatz-like rule. In Michel's notation,
let $g(x)$ be the configuration in state A next to a block of $x$ ones on an otherwise
blank tape. Then
\begin{equation}\label{eq:michel}
g(3k)\to g(5k+6),\qquad g(3k+1)\to g(5k+9),\qquad g(3k+2)\to\text{halt},
\end{equation}
and from $g(0)$ the orbit $0,6,16,34,\dots,12284$ halts after 15 applications
\cite{michelbeh,wikiBB5}.

The bbchallenge paper calls the list of state--symbol pairs visited from the blank tape
the \emph{transcript} of a machine, and uses back-to-back repeats in the transcript to
decide machines that loop forever~\cite[Section~4.3.1]{bbchallenge}. The champion halts,
so its transcript never repeats with a fixed period. But it does repeat with
\emph{growing} counts, and these counts follow~\eqref{eq:michel}.

\paragraph{Contributions.} This paper describes the champion's transcript exactly.
\begin{enumerate}
  \item A closed-form grammar for the whole transcript of $47{,}176{,}870$ symbols
        (Theorem~\ref{thm:main}). Each layer is a short word with explicit exponents.
  \item A Coq proof, with no axioms, that the transcript of the machine from the blank
        tape is exactly this grammar (Section~\ref{sec:coq}). As far as we know, earlier
        machine-checked proofs establish the champion's halting by running it, not by
        its structure~\cite{bbchallenge,coqbb5}.
  \item A proof in Coq that our layer map is the Buro--Michel map~\eqref{eq:michel}
        under the change of variable $a=x+2$ (Proposition~\ref{prop:michel}).
  \item Exact formulas for the steps spent in each layer, which add up to
        $47{,}176{,}870$, and a short account of why the final tape holds $4{,}098$ ones
        (Section~\ref{sec:counts}).
\end{enumerate}
The Collatz-like rule itself is not new. What is new is the transcript-level
description and its formal proof.

% ------------------------------------------------------------------
\section{Preliminaries}

\paragraph{The machine.} The transition table is
\begin{center}
\begin{tabular}{cccccc}
\toprule
 & A & B & C & D & E\\
\midrule
read 0 & 1RB & 1RC & 1RD & 1LA & 1RZ\\
read 1 & 1LC & 1RB & 0LE & 1LD & 0LA\\
\bottomrule
\end{tabular}
\end{center}
The machine starts in state A on a blank tape. We write the halting transition as
$\mathrm{E0}\to\mathrm{1RZ}$; bbchallenge leaves E0 undefined~\cite{bbchallenge}.
Both conventions give $47{,}176{,}870$ steps, and in ours the last step read E0.

\paragraph{Configurations.} We write $\Z\,u\,\hd{s}{x}\,v\,\Z$ for the tape $u\,x\,v$
surrounded by blanks, with the head in state $s$ on the cell holding $x$. Powers denote
repetition, e.g.\ $1^a$ or $(001)^M$.

\paragraph{Transcripts as digits.} The champion has exactly ten state--symbol pairs, so
we write each one as a decimal digit:
\[
d(s,x) = 2\cdot\mathrm{idx}(s) + x,\qquad \mathrm{idx}(\mathrm{A})=0,\dots,
\mathrm{idx}(\mathrm{E})=4.
\]
So A0\,=\,0, A1\,=\,1, B0\,=\,2, \dots, E1\,=\,9. The \emph{transcript} $T$ is the
string of digits $d(s_i,x_i)$ for the steps $i=0,1,\dots,47{,}176{,}869$. This is the
bbchallenge transcript~\cite{bbchallenge} written in digits; the digits only make the
patterns easier to read.

\paragraph{The first 100 digits.} The transcript of the champion begins as follows.
Spaces and line breaks are added only to show where the words of the next section
start; the digits themselves run on without breaks.
\begin{center}
\begin{minipage}{0.8\linewidth}
\begin{verbatim}
0246
159  02  47777603333332  46
159159159  02  4777760333332
47777777776033333333332
4777777777777776033333333333 ...
\end{verbatim}
\end{minipage}
\end{center}
The first line takes the blank tape to the first layer. The second line is the whole
first layer: one $159$, the opening $02$, one rung and the close $46$. From the third
line on, $159$ repeats three times, and each rung is longer than the one before: the
runs of $7$s grow $4, 9, 14$ and the runs of $3$s grow $5, 10, \dots$

% ------------------------------------------------------------------
\section{The transcript grammar}

\subsection{Words}

We use the following words over the digits $0$--$9$.
\begin{center}
\begin{tabular}{lll}
\toprule
Name & Word & Meaning\\
\midrule
$\Omega$ & $159$ & one oscillation A1, C1, E1\\
$R(n)$ & $4\,7^{n+2}\,6\,0\,3^{n+3}\,2$ & internal rung\\
$F(n)$ & $4\,7^{n+2}\,6\,0\,3^{n+4}\,2$ & final rung\\
$G_1$ & $02$ & Gate~1\\
$G_2$ & $1$ & Gate~2\\
$E$ & $46$ & close of a layer\\
$H$ & $158$ & halt\\
\bottomrule
\end{tabular}
\end{center}
In the notation of the earlier versions, a rung $4\,7^{A}\,6\,0\,3^{B}\,2$ has $B=A+1$
when internal and $B=A+2$ when final. We write $K=n+2$ for the length $A$ of the final
rung.

\subsection{Layers}

A layer starts at the configuration
\[
S(a) \;=\; \Z\,1^{a}\,\hd{A}{1}\,1\,\Z ,
\]
with $a\ge 1$ ones to the left of the head. Write $a+1 = 3M + r$ with $0\le r\le 2$.
The layer word $\W(a)$ and the next start $N(a)$ are:
\begin{center}
\renewcommand{\arraystretch}{1.3}
\begin{tabular}{clc}
\toprule
$r$ & $\W(a)$ & $N(a)$\\
\midrule
$0$ & $\Omega^{M}\; G_1\; R(2)\,R(7)\cdots R(5M-8)\; F(5M-3)\; E$ & $5M+3$\\
$1$ & $\Omega^{M}\; G_2\; R(0)\,R(5)\cdots R(5M-5)\; F(5M)\; E$ & $5M+6$\\
$2$ & $\Omega^{M}\; H$ & halt\\
\bottomrule
\end{tabular}
\end{center}
In words: the layer repeats $\Omega$ exactly $M$ times; then the remainder $r$ picks
the gate; then the rungs grow by $5$ each time; a final rung closes the ladder; and
\texttt{46} leads into the next layer. A Gate~1 layer has $M$ rungs and a Gate~2 layer
has $M+1$ rungs, counting the final one. The final rung has
$K = 5M-1$ (Gate~1) or $K = 5M+2$ (Gate~2), and in both cases $N(a) = K+4$.

\begin{theorem}[Transcript grammar]\label{thm:main}
Let $a_0 = 2$ and $a_{i+1} = N(a_i)$. Then $a_{14}+1 \equiv 2 \pmod 3$, and the
transcript of the champion from the blank tape is exactly
\[
T \;=\; 0246\;\W(a_0)\,\W(a_1)\cdots\W(a_{14}).
\]
\end{theorem}

\begin{proof}
Proved in Coq as \coq{bb5\_decimal\_trace}; see Section~\ref{sec:coq}. The prefix
\texttt{0246} takes the blank tape to $S(2)$. Lemmas~\ref{lem:omega}--\ref{lem:bridge}
below give each layer.
\end{proof}

Table~\ref{tab:layers} lists the 15 layers. The first two layers read
\begin{align*}
\W(2) &= 159\;\,02\;\,4\,7^{4}\,6\,0\,3^{6}\,2\;\,46,\\
\W(8) &= (159)^3\;\,02\;\,4\,7^{4}\,6\,0\,3^{5}\,2\;\;4\,7^{9}\,6\,0\,3^{10}\,2\;\;
          4\,7^{14}\,6\,0\,3^{16}\,2\;\,46 .
\end{align*}
The prefix \texttt{0246} is itself $G_1$ followed by $E$: it is a layer with no
oscillations and no rungs.

\begin{table}[t]
\centering
\caption{The 15 layers of the transcript. $a$: ones left of the head at the layer
start; $x=a-2$: Michel's value; $M$: number of $\Omega$; $K$: length $A$ of the final
rung; steps: length of $\W(a)$. The four steps of \texttt{0246} come first; the total
is $47{,}176{,}870$.}
\label{tab:layers}
\small
\begin{tabular}{rrrrcrrr}
\toprule
Layer & $a$ & $x$ & $M$ & Word & Rungs & $K$ & Steps\\
\midrule
1  & 2     & 0     & 1    & Gate 1 & 1    & 4     & 21\\
2  & 8     & 6     & 3    & Gate 1 & 3    & 14    & 83\\
3  & 18    & 16    & 6    & Gate 2 & 7    & 32    & 295\\
4  & 36    & 34    & 12   & Gate 2 & 13   & 62    & 937\\
5  & 66    & 64    & 22   & Gate 2 & 23   & 112   & 2{,}807\\
6  & 116   & 114   & 39   & Gate 1 & 39   & 194   & 8{,}039\\
7  & 198   & 196   & 66   & Gate 2 & 67   & 332   & 22{,}915\\
8  & 336   & 334   & 112  & Gate 2 & 113  & 562   & 64{,}637\\
9  & 566   & 564   & 189  & Gate 1 & 189  & 944   & 180{,}689\\
10 & 948   & 946   & 316  & Gate 2 & 317  & 1582  & 504{,}665\\
11 & 1586  & 1584  & 529  & Gate 1 & 529  & 2644  & 1{,}405{,}029\\
12 & 2648  & 2646  & 883  & Gate 1 & 883  & 4414  & 3{,}908{,}163\\
13 & 4418  & 4416  & 1473 & Gate 1 & 1473 & 7364  & 10{,}864{,}853\\
14 & 7368  & 7366  & 2456 & Gate 2 & 2457 & 12282 & 30{,}201{,}445\\
15 & 12286 & 12284 & 4095 & Halt   & --   & --    & 12{,}288\\
\bottomrule
\end{tabular}
\end{table}

% ------------------------------------------------------------------
\section{Tape lemmas}\label{sec:lemmas}

Each word of the grammar comes from a simple change on the tape. Each lemma below is
proved in Coq for all values of its parameters; the Coq names are given in brackets.
The intermediate configurations are close to those used in the analyses of Buro and
Michel~\cite{buro,michel93}; we state them because the grammar is read directly from
them.

\begin{lemma}[$\Omega$ block; \coq{omega\_block\_tr}]\label{lem:omega}
For $M\ge1$ and any symbol $y$,
\[
\Z\,u\,y\,1^{3M-1}\,\hd{A}{1}\,v\,\Z
\;\xrightarrow{\;(159)^M\;}\;
\Z\,u\,\hd{A}{y}\,(001)^{M}\,v\,\Z .
\]
\end{lemma}
Each $\Omega$ moves the head three cells to the left and turns $111$ into $001$
(A1 writes 1, C1 and E1 write 0). So $M$ oscillations remove $2M$ ones.

\begin{lemma}[Rung; \coq{rung\_tr}]\label{lem:rung}
For $n\ge0$,
\[
\Z\,1^{n}\,\hd{C}{0}\,1\,0\,0\,w
\;\xrightarrow{\;R(n)\;}\;
\Z\,1^{n+5}\,\hd{C}{0}\,w ,
\]
in $2n+9$ steps, where the head ends on the last $0$ of $100$.
\end{lemma}
The rung runs C0, then D sweeps left over $n+2$ ones (the $7$s), D0 and A0 add two ones
at the left end (the $6\,0$), and B sweeps right over $n+3$ ones (the $3$s) before B0
(the $2$). One unit $100$ is filled, and two cells are added on the left: the block
grows by $5$. Repeating the lemma gives a ladder $R(n)\,R(n+5)\cdots$
(\coq{ladder\_tr}).

\begin{lemma}[Final rung; \coq{terminal\_rung\_tr}]\label{lem:final}
For $n\ge0$,
\[
\Z\,1^{n}\,\hd{C}{0}\,1\,1\,\Z
\;\xrightarrow{\;F(n)\;}\;
\Z\,1^{n+6}\,\hd{C}{0}\,\Z ,
\]
in $2n+10$ steps.
\end{lemma}
The extra $3$ in the B sweep is the trailing pair $11$ on the right, which is the
leftover of the start configuration $S(a)$.

\begin{lemma}[Gates, halt and close]\label{lem:bridge}
After the $\Omega$ block of $S(a)$, with $a+1 = 3M+r$:
\begin{itemize}
\item $r=0$: $\Z\,\hd{A}{0}\,(001)^M\,1\,\Z \xrightarrow{\;02\;}
      \Z\,1\,1\,\hd{C}{0}\,(100)^{M-1}\,1\,1\,\Z$;
\item $r=1$: $\Z\,\hd{A}{1}\,(001)^M\,1\,\Z \xrightarrow{\;1\;}
      \Z\,\hd{C}{0}\,(100)^{M}\,1\,1\,\Z$;
\item $r=2$: $\Z\,1\,\hd{A}{1}\,(001)^M\,1\,\Z \xrightarrow{\;158\;}$ halt.
\end{itemize}
After the final rung, $\Z\,1^{m}\,\hd{C}{0}\,\Z \xrightarrow{\;46\;} S(m)$.
\end{lemma}
In Coq these are \coq{layer\_gate1\_tr}, \coq{layer\_gate2\_tr},
\coq{layer\_halt\_tr} and \coq{bridge\_46}. The gate is decided by what is left when the $\Omega$ block reaches the left end of
the block of ones: nothing ($r=0$), one cell ($r=1$) or two cells ($r=2$). With two
cells left, E reads a $0$ and the machine halts.

Combining the lemmas: in a Gate~1 layer the ladder starts at $n=2$ and has $M-1$
internal rungs, so the final rung has $n = 5M-3$ and the next start is
$a' = n+6 = 5M+3$. In a Gate~2 layer the ladder starts at $n=0$ and has $M$ internal
rungs, so $n = 5M$ and $a' = 5M+6$. This is the table for $N(a)$.

% ------------------------------------------------------------------
\section{Relation to the Buro--Michel map}

\begin{proposition}[\coq{layer\_map\_is\_michel}]\label{prop:michel}
For every $x\ge0$, the layer map $N$ at $a=x+2$ is the map~\eqref{eq:michel} shifted
by $2$: if $x=3k$ then $N(x+2) = (5k+6)+2$; if $x=3k+1$ then $N(x+2) = (5k+9)+2$; if
$x=3k+2$ the layer halts.
\end{proposition}

So the layer starts $a_i$ are Michel's orbit $x_i$ plus $2$ (Table~\ref{tab:layers}).
Since $N(a)=K+4$, Michel's next value is $x' = K+2$. Two consequences, which restate
the known rule in the terms of the grammar:
\begin{corollary}[\coq{gate\_law}, \coq{recurrence\_8\_1}]\label{cor:gate}
Let $K$ be the final-rung length of a layer. The next layer has
\[
M' = \left\lfloor \tfrac{K+2}{3} \right\rfloor + 1,
\]
and it is a Gate~1 layer if $(K+2)\bmod 3 = 0$, a Gate~2 layer if $(K+2)\bmod 3 = 1$,
and the halting layer if $(K+2)\bmod 3 = 2$.
\end{corollary}
For the champion, the remainder $2$ first appears after layer~14, where $K+2=12284$.

% ------------------------------------------------------------------
\section{Steps and ones}\label{sec:counts}

\paragraph{Steps.} The step counts in Lemmas~\ref{lem:omega}--\ref{lem:bridge} give the
length of each layer word:
\[
|\W(a)| =
\begin{cases}
5M^2 + 11M + 5 & \text{Gate 1},\\
5M^2 + 17M + 13 & \text{Gate 2},\\
3M + 3 & \text{halt}.
\end{cases}
\]
With the four steps of \texttt{0246}, the 15 layers add up to exactly $47{,}176{,}870$
steps (Table~\ref{tab:layers}). The formulas are sums of the proved step counts; the
total is checked by exact integer arithmetic. It is not yet stated inside Coq, because
Coq's unary natural numbers make a number of this size impractical; a version with
binary numbers is future work. The total agrees with the stage counts of
Michel~\cite{michelbeh}, which use different stage boundaries.

\paragraph{Ones.} The start $S(a)$ holds $a+2$ ones. Writing $a+1 = 3M+r$ and
$g = r+1$, this is $3M + g$, and after the $\Omega$ block there are $M+g$ ones
(Lemma~\ref{lem:omega}). In the halting layer, $M = 4095$ and $g = 3$, leaving
$4095+3 = 4098$ ones. The final word $158$ does not change this: C1 erases a one and
E0 writes one. This gives $\Sigma(5) = 4098$ for the champion.

% ------------------------------------------------------------------
\section{The Coq proof}\label{sec:coq}

The development is the file \coq{coq/BB5Omega.v} in the repository (about 840 lines,
Coq~8.18). It uses no axioms: \coq{Print Assumptions} on the main theorems reports
that they are closed under the global context.

\paragraph{Model.} A configuration is a state, the list of cells to the left of the
head (nearest first), the cell under the head, and the list of cells to the right. An
empty list stands for blanks. \coq{step} applies the transition table, and
\coq{trace n c} lists the digits $2\cdot\mathrm{idx}(s)+x$ of the first $n$ steps from
$c$. The key fact used everywhere is
\[
\coq{trace}\,(p+q)\,c \;=\; \coq{trace}\,p\,c \;\cdot\; \coq{trace}\,q\,(\coq{steps}\,p\,c).
\]

\paragraph{Proof structure.} Each tape lemma is proved twice: once for the
configuration it reaches and once for the digits it writes. Sweeps and blocks of any
length are handled by induction. The layer lemmas are put together into
\coq{layer\_correct\_tr}, which covers every $a\ge1$. A short recursion then runs the
layer map from $a=2$; Coq computes that it halts after 15 layers
(\coq{fifteen\_layers}) and the sequence of $M$ values (\coq{M\_sequence}).

\paragraph{Main statement.} Here \coq{grammar 20 2} is the concatenation of the layer
words starting from $a=2$ (20 is an upper bound on the number of layers):
\begin{lstlisting}
Theorem bb5_decimal_trace :
  exists n, state (steps n (A, [], false, [])) = Z
         /\ trace n (A, [], false, []) = [0; 2; 4; 6] ++ grammar 20 2.
\end{lstlisting}
As an extra check, we computed the first five layers of \coq{grammar} inside Coq
($4{,}147$ digits) and compared them with an independent simulation; they agree. A
separate Python script in the repository, \coq{verify\_gate\_law.py}, runs all
$47{,}176{,}870$ steps and checks the layer data of Table~\ref{tab:layers} from the
transcript.

\paragraph{What is not in Coq.} The total step count and the final count of ones are
derived from the proved lemmas by arithmetic, as in Section~\ref{sec:counts}, but are
not stated as Coq theorems.

% ------------------------------------------------------------------
\section{A smaller case: BB(4)}

The same lens applied to the 4-state winner
\texttt{1RB1LB\_1LA0LC\_1RZ1LD\_1RD0RA} ($107$ steps, $13$ ones; digits $0$--$7$)
gives
\[
0212\;(03567)^3\;0212\,1\,3567\;(03567)^4\;0212\,1\,357\,1212\;(03567)^4\;
0212\,1\,357\,12\;(03567)^2\;0212\,1\,34 .
\]
There is a repeated word ($03567$), a gate-like opening ($0212$) and a halt by a change
of one digit at the end of the repeated word ($35\to34$, as $159\to158$ for BB(5)). But
the counts $3,4,4,2$ do not grow, and there is no layer map. This is only an
observation on one small run; we do not claim a general pattern.

% ------------------------------------------------------------------
\section{Discussion}

\paragraph{Growing repeats.} The bbchallenge decider for loops looks for a transcript
block that repeats back-to-back with a fixed length~\cite[Algorithm~1]{bbchallenge}. The
champion's transcript is made of blocks that repeat a growing number of times, and the
counts follow a Collatz-like map. A natural question is whether such grammars can be
found automatically: detect growing repeats in a transcript, fit a layer map from a few
layers, and then iterate the map instead of the machine. For BB(5) the first layers can
mislead: $M = 1,3,6,12,22$ looks like doubling before the ratio settles near $5/3$. So
several layers, and every remainder class, must be seen before a map is trusted. Even a
correct map does not settle halting in general: for machines such as
Antihydra~\cite{bbchallenge} the question becomes a hard Collatz-like problem.

\paragraph{Formal proofs of champions.} The Coq proof of $S(5)$ needs only that the
champion halts within $47{,}176{,}870$ steps, which it checks by running
it~\cite{bbchallenge}. A structural proof, as here, says why. Proofs of this kind may
be useful for larger champions, whose runs are far too long to execute.

\paragraph{Future work.} (1) State the step count and the count of ones in Coq, using
binary numbers. (2) Try transcript grammars on other Collatz-like machines, including
BB(6) champions. (3) Extract layer maps from transcripts automatically.

% ------------------------------------------------------------------
\section{Conclusion}

The transcript of the BB(5) champion has an exact grammar: \texttt{0246} followed by 15
layers of the form $(159)^M$, a gate, a ladder of rungs that grow by five, a final rung
and \texttt{46}, with a last layer ending in \texttt{158}. The grammar is the
Buro--Michel map seen at the level of single steps. We proved in Coq, without axioms,
that the machine's transcript is exactly this grammar.

\section*{Acknowledgements}
Parts of the analysis and of
the Coq development were carried out with the help of an AI assistant (Claude,
Anthropic); the author checked all results, and the Coq proofs are checked by the Coq
kernel.

\begin{thebibliography}{99}

\bibitem{bbchallenge}
The bbchallenge Collaboration: J.~Blanchard, D.~Briggs, K.~Deka, N.~Fenner,
Y.~Forster, G.~Georgiev (Skelet), M.~L.~House, R.~Hunter, Iijil, M.~K\k{a}dzio\l{}ka,
P.~Kropitz, S.~Ligocki, mxdys, M.~Na\'sciszewski, savask, T.~St\'erin, C.~Xu, J.~Yuen,
T.~Zimmermann.
\newblock Determination of the fifth Busy Beaver value.
\newblock arXiv:2509.12337, 2025.

\bibitem{coqbb5}
mxdys.
\newblock Coq-BB5: a Coq proof that $S(5) = 47{,}176{,}870$.
\newblock \url{https://github.com/ccz181078/Coq-BB5}, 2024.

\bibitem{marxen}
H.~Marxen and J.~Buntrock.
\newblock Attacking the Busy Beaver 5.
\newblock \emph{Bulletin of the EATCS}, 40:247--251, 1990.

\bibitem{buro}
M.~Buro.
\newblock Ein Beitrag zur Bestimmung von Rados $\Sigma(5)$ oder Wie f\"angt man fleissige Biber?
\newblock Schriften zur Informatik und angewandten Mathematik, Report No.~146, RWTH Aachen, 1990.

\bibitem{michel93}
P.~Michel.
\newblock Busy beaver competition and Collatz-like problems.
\newblock \emph{Archive for Mathematical Logic}, 32:351--367, 1993.

\bibitem{michelsurvey}
P.~Michel.
\newblock The Busy Beaver competition: a historical survey.
\newblock arXiv:0906.3749, 2022.

\bibitem{michelbeh}
P.~Michel.
\newblock Behavior of busy beavers.
\newblock \url{https://bbchallenge.org/~pascal.michel/beh}.

\bibitem{wikiBB5}
BusyBeaverWiki.
\newblock BB(5) champion.
\newblock \url{https://wiki.bbchallenge.org/wiki/BB(5)_champion}.

\bibitem{aaronson}
S.~Aaronson.
\newblock The Busy Beaver frontier.
\newblock \emph{SIGACT News}, 51(3):32--54, 2020.

\bibitem{rado}
T.~Rad\'o.
\newblock On non-computable functions.
\newblock \emph{Bell System Technical Journal}, 41(3):877--884, 1962.

\end{thebibliography}

\end{document}
